\section{Running compression at a prescribed target roof}
\label{sec:v97-running}
The next estimate controls a maximum over the sampled time before
projection and before the collision limit. Its norm is essential
height in the target roof, not an integral over starting roofs.
The sampled interval is deliberately kept in the original exact event:
no singular-volume theorem on the enlarged comparison domain is assumed.

Fix $a,r>0$, $0<q\le1/2$, and put $d=2qr$. Using the original
stopped lifetimes, define disjoint target classes
\begin{align}
 P_-&=\{\tau^->2d\},\notag\\
 P_+&=\{\tau^-\le2d,\ z\ge a,\ \tau^+>d\},\qquad P=P_-\cup P_+.
                                      \label{eq:v97-protected}
\end{align}
Choose sign $e=-1$ on $P_-$ and $e=1$ on $P_+$. The chosen interval
of length $d$ lies entirely in the same original exact-event segment.
Set
\[
 \rho=\min\{a,\sqrt{2\lambda_\circ d}\},\qquad
 v^\rho(u)\,du=F_\#(I_mz^{-1}\one_{\{z\ge\rho\}}\bar\nu).
\]
On $P_+$ radial growth gives $z(\Phi_sx)\ge a$; on $P_-$ it gives
$z(\Phi_{-s}x)^2\ge2\lambda_\circ(\tau^-(x)-s)>2\lambda_\circ d$
for $0\le s\le d$. Thus the same positive cutoff $\rho$ works on
both classes, including every small internal incidence.

For $D_0>1$ define the original weighted running-compression source
\begin{equation}\label{eq:v97-running-source}
 C_{a,q,r,D_0}^\varepsilon(t)=
 \int_{F=t}W^\varepsilon\one_P
 \one_{\{\sup_{0\le s\le d}A(x)/A(\Phi_{es}x)>D_0\}}\,d\mu_t(x).
\end{equation}
The supremum can be taken over rational times and endpoints.

\begin{theorem}[A complete target-height bound for running loss]
\label{thm:v97-running}
At every finite count and almost every roof,
\begin{equation}\label{eq:v97-running-finite}
 C_{a,q,r,D_0}^\varepsilon(t)
 \le\frac C{1-D_0^{-1}}\int_{t-d}^{t+d}v^\rho(u)\,du.
\end{equation}
Consequently, for all previously fixed $a,q,r,D_0$,
\begin{equation}\label{eq:v97-running-height}
 \limsup_m\sup_{R,n,k}m^2\|C_{a,q,r,D_0}^\varepsilon\|_\infty
 \le\frac{Cqr}{1-D_0^{-1}}.
\end{equation}
The leading constant is independent of $a$ and $\varepsilon$.
Neither $a=0$ nor a collision-dependent cutoff is asserted.
\end{theorem}
\begin{proof}
For a target counted in \eqref{eq:v97-running-source}, continuity
provides a sampled time at which $A_s<A_0/D_0$. Hence
\[
 (1-D_0^{-1})A_0\le\int_0^d
                  |\partial_u A(\Phi_{es}x)|\,ds.
\]
The factor on the left pays the original target density, not the
terminal density. Multiply by $W^\varepsilon(x)\le1$ and use
$|\partial_uA|\le CAI_m/z$. At every sampled roof the Borel
segment bijection counts the selected states once, inside the same
exact event. The lower bound for $z$ therefore identifies their
sum as a submeasure of $v^\rho$ at that roof. The negative and
positive target-sign classes use disjoint halves of $(t-d,t+d)$.
This proves \eqref{eq:v97-running-finite} by Tonelli, without
removing the density ratio or taking a positive part after a false
change of source. Apply the positive-cutoff volume estimate
\eqref{eq:v94-cutoff} to the fixed interval of length $2d$ to obtain
\eqref{eq:v97-running-height}. Its finite-count remainder may depend
on $\rho$; the leading constant after the stated limit does not.
\end{proof}

\begin{corollary}[Low capped occupation on the whole protected class]
\label{cor:v97-low-clock}
For $H=\Pi$,
\begin{align}
 R_{P\cap\{c_r<q\}}^\varepsilon(t)
       &\le C\int_{t-d}^{t+d}v^\rho(u)\,du,
                                      \label{eq:v97-low-clock-finite}\\
 \mathcal H_M(R_{P\cap\{c_r<q\}}^\varepsilon)&\le Cqr.
                                      \label{eq:v97-low-clock-height}
\end{align}
\end{corollary}
\begin{proof}
If $A_s/A_0\ge1/2$ throughout the chosen original interval, that
interval alone contributes at least $d/(2r)=q$ to $c_r$. Therefore
$c_r<q$ implies that the running ratio exceeds two. Apply
Theorem~\ref{thm:v97-running} with $D_0=2$.
\end{proof}

In particular the running loss is paid even if a later part of the
physical comparison encounters many section decisions or fails the
old worst-time test. The proof samples only the displayed initial
same-event interval; it does not estimate optical variation accumulated
at states carrying a different exact occupation label.
